\chapter{总结与展望}

\section{全文总结}

本研究面向中性粒细胞这一终末分化、处于高动态应激状态的免疫细胞，围绕第一章提出的核心问题展开：在原代中性粒细胞高通量功能筛选受限的条件下，如何基于深度学习必需性预测模型识别免疫情境中的关键功能依赖蛋白。围绕这一问题，本文按照“模型预测—候选筛选—机制验证”的总体思路，构建了跨层级蛋白必需性预测框架，并结合结构与功能证据对核心候选蛋白进行了分析与验证。本文在生物学解释与实验设计层面进一步将免疫情境具体化为急性炎症应答中以氧化爆发增强和颗粒/溶酶体功能负荷升高为特征的活化中性粒细胞状态，即氧化爆发-颗粒/溶酶体高负荷活化状态。整体而言，本研究在方法构建、候选筛选与机制解释三个层面均取得了较为系统的进展。

围绕上述研究目标，本文的主要研究工作与结果可概括为以下三个方面。

第一，构建了面向免疫情境依赖的深度学习预测模型。本文采用 ProtT5 预训练蛋白语言模型提取序列嵌入表示，在此基础上构建多种分类头并开展系统消融比较。结果显示，相较于仅依赖池化表示的基线模型，注意力增强结构能够更有效捕捉与必需性判别相关的关键残基信息。在本研究任务设定下，单头注意力结构在性能与稳定性之间取得更优平衡，最终模型在独立测试集上达到 AUROC 0.9272 和 AUPRC 0.8245，说明所构建框架能够在类别不均衡场景中保持较强的识别能力。

第二，建立了 Human-level 与 Immune-level 双层评分体系，并据此解析免疫情境候选必需蛋白。本文通过构建跨层级 protein essential score（PES）并引入差异指标 $\Delta$PES，对“基础生存依赖”与“免疫情境条件性依赖”进行区分。基于该框架识别出的高置信候选蛋白并未随机分布，而是在功能与性质层面呈现出清晰聚类。进一步理化特征分析表明，免疫情境候选必需蛋白整体体现出正电荷偏移、分子量轻量化与疏水核心强化等趋势，这些变化与活化中性粒细胞在酸性吞噬体环境、强氧化应激及高膜周转负荷下的功能约束方向一致。

第三，围绕核心候选蛋白形成了较为清晰的机制解释链条。结合注意力图谱与文献证据，本文重点关注 ATP6V1B2 与 PLBD1（LyLAP）。结果提示，ATP6V1B2 在 Human-level 上得分较低（PES=0.007），而在 Immune-level 上显著升高（PES=0.952），其作用更可能体现为氧化爆发-颗粒/溶酶体高负荷活化状态下的酸化动力学调控节点，而非一般意义上的基础生存核心因子。对于 PLBD1，本文结果支持其在溶酶体膜稳态维护中的潜在关键作用，即通过处理疏水肽段降低膜扰动风险。上述发现共同说明，模型不仅能够输出“是否重要”的统计判别，还可为“在何种生理背景下重要”提供机制导向线索。

进一步看，本研究的创新点与研究意义主要体现在以下三个方面。

（1）在问题定义层面，本文将蛋白必需性从单一静态标签扩展为跨层级、情境依赖的连续刻画框架。相较于仅报告核心必需与非必需二元分类的传统做法，PES 与 $\Delta$PES 提供了更细粒度的比较维度，有助于区分“基础维持必需”与“免疫情境相关依赖”。

（2）在方法路径层面，本文验证了“序列驱动 + 轻量注意力判别头”在免疫情境依赖预测任务中的可行性。该路径在不依赖昂贵结构实验输入的前提下，仍可实现较高判别性能，为原代免疫细胞中难以直接开展高通量编辑筛选的问题提供了可扩展的计算补充方案。

（3）在应用价值层面，本文筛选得到的候选集合及其机制线索，为后续干预靶点优选提供了依据。尤其是兼具较高 Immune-level 依赖性与较低 Human-level 依赖性的候选蛋白，理论上更具“选择性调控”潜力，可能在降低全身副作用的同时实现对炎症相关过程的有效干预。

\section{研究局限性分析}

尽管本文形成了较完整的计算筛选与机制解释流程，但仍存在以下限制，需要在后续研究中进一步完善。

在数据层面，训练标签主要来源于既有大规模筛选数据库与间接映射流程，数据本身受实验体系、细胞背景与统计校正策略影响。即使经过标准化处理，不同来源数据之间仍可能存在不可忽略的批次效应与语境偏差。此外，本研究的 Immune-level 主要反映免疫效应执行相关高应激功能情境中的依赖特征，尚不能完全覆盖原代中性粒细胞在真实炎症微环境中的状态连续变化。尽管本文在实验验证中进一步聚焦于氧化爆发-颗粒/溶酶体高负荷活化状态，但当前结论更适宜理解为“在既定数据语境下具有较高可信度的推断”，其外推范围仍需在更广泛免疫状态中进一步检验。

在模型层面，本文以蛋白序列为核心输入，虽具有通用性与可扩展性，但也天然忽略了部分与必需性密切相关的动态信息，包括表达时序变化、翻译后修饰、亚细胞定位动态迁移以及蛋白互作网络重构等。当前模型对这些因素的捕捉主要依赖预训练表示中的隐式编码，尚未实现显式建模。此外，本文采用的判别头结构以性能与可训练稳定性为优先，尚未系统探索更复杂的不确定性估计与因果结构约束，对极端少样本情境的鲁棒性仍有提升空间。

在验证层面，本研究的实验支撑主要来自文献一致性、结构解析线索与部分替代体系验证，尚未在原代中性粒细胞中对全部核心候选蛋白开展统一、同尺度的功能扰动实验。考虑到中性粒细胞寿命短、体外维持困难及基因操作窗口受限，当前“计算预测到功能因果”的证据层级仍然有限。后续若补充更直接的原代扰动实验，可进一步增强结论的因果解释力度与外部效度。

\section{未来工作展望}

围绕上述局限，后续研究可从以下方向展开。

首先，可在现有序列建模框架基础上，引入中性粒细胞不同激活阶段的动态转录组、蛋白质组及空间组学信息，构建显式情境变量约束。通过将细胞状态嵌入与序列表示联合学习，可进一步提高模型对“同一蛋白在不同免疫状态下功能权重变化”的感知能力，从而增强跨条件外推性能。

其次，在实验策略上，建议优先围绕高 $\Delta$PES 候选建立小规模、高质量验证体系。可结合快速蛋白降解技术，如蛋白水解靶向嵌合体（proteolysis targeting chimera, PROTAC）和 dTAG 靶向降解标签系统（degradation TAG system, dTAG），并配合时间分辨表型评估，在原代中性粒细胞中直接观察扰动后酸化动力学、膜完整性、吞噬体成熟与 NETs 相关表型变化。该策略有助于将模型输出由“候选排序”进一步推进到“具备因果证据支持的功能节点”。

再次，对于 ATP6V1B2、PLBD1 等已显示免疫情境候选必需特征的蛋白，后续可进一步构建“有效性-安全性”联合评价框架，综合 Human-level 与 Immune-level 评分、组织表达谱、潜在脱靶风险及可成药性指标，形成分层优先级列表。在此基础上推进小分子或生物调控剂开发，有望为慢性炎症、自身免疫相关病理过程及血栓相关疾病提供更具选择性的干预思路。

最后，本文提出的跨层级建模与候选筛选流程并不局限于中性粒细胞。未来可迁移至单核/巨噬细胞、T细胞或组织驻留免疫细胞等体系，形成可复用的免疫细胞依赖性研究范式。通过在多细胞类型中的横向验证与对照，进一步识别细胞共性高依赖模块与细胞情境特异依赖模块，为免疫系统层面的精细化靶点发现提供方法基础。

综上，本文在中性粒细胞情境下完成了从模型构建到候选机制解析的初步研究流程，证明了序列驱动深度学习在免疫情境候选必需蛋白识别中的可行性。随着多模态数据融合与原代验证体系的持续推进，该研究路线有望进一步提升在机制阐释与临床转化中的应用价值。
